\subsection{state}\label{state}
Die Struktur \verb|state_t| ist die zentrale Datenstruktur des Systems: sie
hält den Tastenzustand, der zwischen allen Eingabegeräten (Abschnitt
\ref{keyboard}) ausgetauscht und über Funk zwischen den Geräten
synchronisiert wird. Jede der bis zu \verb|MAX_BUTTON_CNT| Tasten besitzt
einen Zustand \verb|key_state_e| -- \verb|OFF|, \verb|BLINKING| oder
\verb|ON| (Listing \ref{key_state_e}) -- sowie eine Beschriftung.

Der zentrale Mechanismus ist \verb|state_propagate()|: ein Tastendruck
schaltet den Zustand zyklisch weiter
($\mathtt{OFF} \to \mathtt{BLINKING} \to \mathtt{ON} \to \mathtt{OFF}$,
also $\mathtt{state} = (\mathtt{state}+1) \bmod \mathtt{STATE\_CNT}$) und
markiert den \verb|state_t| als \emph{dirty}.

\begin{listing}
\begin{minted}{C}
typedef enum __attribute__((packed)) {
    OFF, BLINKING, ON, STATE_CNT
} key_state_e;
\end{minted}
\caption{Tastenzustand \texttt{key\_state\_e}}
\label{key_state_e}
\end{listing}

\subsubsection*{Zustand}
Listing \ref{state_t} zeigt die Struktur (Größe $40$~Byte). \verb|first| und
\verb|cnt| legen das aktive Fenster der Tasten fest, \verb|dirty| markiert
eine Änderung seit der letzten Synchronisation, \verb|id| ist ein bei jedem
Senden mitgeführter Zähler und \verb|clabel| (Abschnitt \ref{common}) trägt
ein Kommando bzw.\ Label. Die Felder \verb|state[]| und \verb|label[]| halten
Zustand und Beschriftung je Taste.

\begin{listing}
\begin{minted}{C}
typedef struct state_s {
    uint8_t  first;                 // erste gueltige Taste
    uint8_t  cnt;                   // Anzahl gueltiger Tasten
    uint8_t  dirty;                 // Aenderungs-Flag
    uint8_t  id;                    // Sende-Zaehler
    clabel_u clabel;                // Kommando/Label (4 Byte)
    key_state_e state[MAX_STATE_CNT]; // Zustand je Taste
    char     label[MAX_STATE_CNT];  // Beschriftung je Taste
} state_t;                          // 2*MAX_BUTTON_CNT + 8 = 40 Byte
\end{minted}
\caption{Definition der \texttt{state\_t} Struktur}
\label{state_t}
\end{listing}

\noindent Über die Option \verb|REDUCED_PAYLOAD| (in \verb|emLibC_options.h|)
existiert alternativ eine bit\-gepackte Variante ($\approx 10$~Byte), die hier
deaktiviert ist.

\subsubsection*{Lebenszyklus und Einzeltasten}
{\sloppy\emergencystretch=3em
\begin{description}
\item[\texttt{state\_init(state)} / \texttt{state\_reset(state)}] Initialisiert
  die Struktur (Standardbeschriftung \verb|0..F|) bzw.\ setzt alle Tasten auf
  \verb|OFF|. Rückgabe: \verb|em_msg| (\verb|EM_ERR| bei ungültiger Struktur,
  sonst \verb|EM_OK|).
\item[\texttt{state\_check(state)}] Prüft die Bereichsfelder auf Gültigkeit.
  Rückgabe: \verb|em_msg| (\verb|EM_OK| bei gültigen Bereichsfeldern, sonst
  \verb|EM_ERR|).
\item[\texttt{state\_set(state, nr, ns)} / \texttt{state\_get(state, nr)}]
  Setzt bzw.\ liest den Zustand einer Taste über ihren Index. Rückgabe:
  \verb|em_msg| (Setzer; \verb|EM_ERR| bei ungültigem Index/Struktur) bzw.\
  \verb|key_state_e| (der Tastenzustand).
\item[\texttt{state\_set\_key\_by\_lbl(\dots)} /
  \texttt{state\_get\_key\_by\_lbl(\dots)}] Dasselbe über die Beschriftung;
  \verb|state_ch2idx()| bildet ein Zeichen auf den Index ab (\verb|int8_t|,
  $-1$ wenn nicht gefunden). Rückgabe: \verb|em_msg| (Setzer) bzw.\
  \verb|key_state_e| (der Tastenzustand, \verb|STATE_CNT| wenn nicht gefunden).
\item[\texttt{state\_propagate(state, idx)} /
  \texttt{state\_propagate\_by\_idx} / \texttt{\dots\_by\_lbl}] Schaltet eine
  Taste zyklisch um einen Schritt weiter und setzt \emph{dirty}. Rückgabe:
  \verb|em_msg| (\verb|EM_ERR| bei ungültiger Struktur, sonst \verb|EM_OK|).
\end{description}
\par}

\subsubsection*{Ganzer Zustand: Übertragen und Kombinieren}
{\sloppy\emergencystretch=3em
\begin{description}
\item[\texttt{state\_set\_state(from, to)} / \texttt{state\_copy(from, to)}]
  Übernehmen den Tasteninhalt von \verb|from| nach \verb|to|. Rückgabe:
  \verb|em_msg| (\verb|EM_ERR| bei ungültiger Struktur, sonst \verb|EM_OK|).
\item[\texttt{state\_set\_u32(state, u32)} / \texttt{state\_get\_u32(state)}]
  Packen bzw.\ entpacken die $16$ Tasten zu je $2$ Bit in einen
  \verb|uint32_t|. Rückgabe: \verb|em_msg| (Setzer) bzw.\ \verb|uint32_t| (die
  gepackten Zustände).
\item[\texttt{state\_diff(ref, state, diff)}] Bestimmt die tasten\-weise
  Differenz (modular, vgl.\ \verb|state_key_diff()|) gegenüber \verb|ref| in
  \verb|diff|. Rückgabe: \verb|em_msg| (das \emph{dirty}\-Flag).
\item[\texttt{state\_add(ref, add)}] Wendet \verb|add| als Anzahl von
  Weiterschaltungen auf \verb|ref| an (\verb|BLINKING| $= +1$, \verb|ON| $=
  +2$). Rückgabe: \verb|em_msg| (das \emph{dirty}\-Flag).
\item[\texttt{state\_merge(in, out)}] Übernimmt geänderte Tasten von \verb|in|
  nach \verb|out| und meldet, ob sich etwas geändert hat. Rückgabe:
  \verb|bool|.
\item[\texttt{state\_is\_same(last, cur)}] Prüft zwei Zustände auf Gleichheit.
  Rückgabe: \verb|em_msg| (\verb|EM_OK| bei Gleichheit).
\end{description}
\par}

\subsubsection*{Bereich, Flags und Ausgabe}
{\sloppy\emergencystretch=3em
\begin{description}
\item[\texttt{state\_set\_first/cnt} / \texttt{state\_get\_first/cnt}] Setzen
  bzw.\ lesen das aktive Tastenfenster. Rückgabe: \verb|em_msg| (Setzer) bzw.\
  \verb|uint8_t| (Getter).
\item[\texttt{state\_get\_dirty} / \texttt{state\_set\_dirty} /
  \texttt{state\_set\_undirty}] Verwalten das Änderungs\-Flag
  (\verb|state_get_dirty| liefert dessen Wert). Rückgabe: \verb|em_msg|
  (\verb|EM_ERR| bei ungültiger Struktur).
\item[\texttt{state\_print(state, title, doLong, cycle)}] Gibt Beschriftung,
  Zustände und die Cycle\-Position (Abschnitt \ref{cycle}) zu Debug\-Zwecken
  aus. Rückgabe: \verb|em_msg| (\verb|EM_ERR| bei ungültiger Struktur, sonst
  \verb|EM_OK|).
\end{description}
\par}
